% Part: computability
% Chapter: recursive-functions
% Section: non-pr-functions

\documentclass[../../../include/open-logic-section]{subfiles}

\begin{document}

\olfileid{cmp}{rec}{npr}
\olsection{دوال غير عودية بدائية}

لا تستنفد الدوال العودية البدائية جميع الدوال القابلة للحساب بحسب
الحدس. وينبغي أن يكون واضحًا حدسيًا أننا نستطيع وضع قائمة بكل الدوال
العودية البدائية الأحادية الحجة $f_0$ و$f_1$ و$f_2$ و…، على نحو
يمكننا معه حساب قيمة $f_x$ عند الدخل $y$ حسابًا فعالًا؛ وبعبارة أخرى،
تكون الدالة $g(x,y)$ المعرفة بـ
\[
g(x,y) = f_x(y)
\]
قابلة للحساب. ولكن الدالة الآتية تكون قابلة للحساب أيضًا:
\begin{eqnarray*}
h(x) & = & g(x,x) + 1 \\
& = & f_x(x) +1.
\end{eqnarray*}
ولكل دالة عودية بدائية $f_i$، تختلف قيمة $h$ وقيمة $f_i$ عند $i$. ومن ثم
تكون $h$ قابلة للحساب ولكنها ليست عودية بدائية؛ ويمكن قول الشيء نفسه
عن $g$. وهذه نسخة «فعالة» من حجة كانتور القطرية.

ويمكن إعطاء أمثلة أصرح لدوال قابلة للحساب وليست عودية بدائية. فليدل
الترميز $g^n(x)$، مثلًا، على $g(g(\dots g(x)))$، وفيه $n$ من ورود
$g$؛ ولنعرف متتالية الدوال $g_0,g_1,\dots$ بـ
\begin{eqnarray*}
g_0(x) & = & x+1 \\
g_{n + 1}(x) & = & g_n^x(x)
\end{eqnarray*}
يمكنك التحقق من أن كل دالة $g_n$ عودية بدائية. وكل دالة تالية تنمو
أسرع كثيرًا من سابقتها؛ فـ$g_1(x)$ تساوي $2x$، و$g_2(x)$ تساوي
$2^x\cdot x$، وتنمو $g_3(x)$ نموًا يقارب برجًا أسيًا يتألف من $x$
نسخ من العدد~$2$. ودالة أكرمان--بيتر هي في جوهرها الدالة $G(x)=g_x(x)$،
ويمكن بيان أنها تنمو أسرع من أي دالة عودية بدائية.

لنعد إلى مسألة تعداد الدوال العودية البدائية. تذكر أننا أسندنا
ترميزات رمزية إلى كل دالة عودية بدائية، ولذلك يكفي تعداد الترميزات.
يمكننا أن نسند عددًا طبيعيًا $\#(F)$ إلى كل ترميز $F$ عوديًا كما يأتي:
\begin{eqnarray*}
\#(0) & = & \langle 0 \rangle \\
\#(S) & = & \langle 1 \rangle \\
\#(\Proj{n}{i}) & = & \langle 2, n, i \rangle \\
\#(\fn{Comp}_{k,l}[H,G_0,\dots,G_{k-1}]) & = & \langle
3,k,l,\#(H),\#(G_0),\dots,\#(G_{k-1}) \rangle \\
\#(\fn{Rec}_l[G,H]) & = & \langle 4, l, \#(G), \#(H) \rangle
\end{eqnarray*}
نستعمل هنا حقيقة أن كل متتالية أعداد يمكن النظر إليها بوصفها عددًا
طبيعيًا، باستعمال رموز القسم السابق. وخلاصة ذلك أن كل ترميز يسند إليه
عدد طبيعي. ولا تقابل بعض المتتاليات، ومن ثم بعض الأعداد، ترميزات
بالطبع؛ لكن يمكننا أن نجعل $f_i$ الدالة العودية البدائية الأحادية
الحجة ذات الترميز الذي رمزه $i$ إذا كان $i$ يرمز ترميزًا كهذا، ونجعلها
الدالة الثابتة $0$ في غير ذلك. والنتيجة النهائية أن لدينا طريقة صريحة
لتعداد الدوال العودية البدائية الأحادية الحجة.

(والواقع أن بعض الدوال، كالدالة الثابتة صفرًا، ستظهر أكثر من مرة في
القائمة. وليس هذا مجرد أثر مصطنع لترميزنا، بل هو أيضًا نتيجة لوجود
أكثر من ترميز للدالة الثابتة صفرًا. وسنرى لاحقًا أنه لا يمكن تجنب هذه
التكرارات على نحو قابل للحساب؛ فمثلًا لا توجد دالة قابلة للحساب تقرر
ما إذا كان ترميز معطى يمثل الدالة الثابتة صفرًا.)

يمكننا الآن أخذ الدالة $g(x,y)$ معطاة بـ$f_x(y)$، حيث تشير $f_x$ إلى
التعداد الذي وصفناه للتو. كيف نعرف أن $g(x,y)$ قابلة للحساب؟ هذا واضح
حدسيًا: فلحساب $g(x,y)$، «نفك» $x$ أولًا وننظر هل هو ترميز لدالة
أحادية الحجة. فإن كان كذلك، حسبنا قيمة تلك الدالة عند الدخل~$y$.

\begin{tagblock}{TMs}
\begin{digress}
قد تكون مقتنعًا من قبل بأنه يمكن، مع بعض العمل، كتابة برنامج، بلغة
جافا أو ++C مثلًا، يفعل ذلك؛ ويمكننا الآن الاحتكام إلى أطروحة
تشرش--تورنغ، التي تقول إن كل ما هو قابل للحساب بحسب الحدس يمكن أن
تحسبه آلة تورنغ.

وثمة بالطبع طريقة أكثر مباشرة لبيان أن $g(x,y)$ قابلة للحساب، هي وصف آلة
تورنغ تحسبها وصفًا صريحًا. وهذا يتجنب، بوجه خاص، أطروحة تشرش--تورنغ
والاحتكام إلى الحدس. وسنكون قد بنينا قريبًا من الأدوات ما يكفي لبيان
أن $g(x,y)$ قابلة للحساب بالاحتكام إلى نموذج حساب يمكن
\emph{محاكاته} على آلة تورنغ، وهو الدوال العودية.
\end{digress}
\end{tagblock}

\end{document}
